% ============================================================
% 双栏中文学术论文模版 — Windows 版
% 中文字体：Windows 系统字体（宋体/黑体/楷体/仿宋）
% 英文字体：Times New Roman / Arial / Courier New（Windows 自带）
% 编译：xelatex main-win.tex
% ============================================================
\documentclass[10pt, twocolumn, a4paper]{article}

% ---- 基础编码与中文支持 ----
\usepackage[UTF8, heading=true, fontset=windows]{ctex}
\ctexset{
  section/format       = \heiti\zihao{-4}\bfseries,
  subsection/format    = \heiti\zihao{5}\bfseries,
  subsubsection/format = \kaishu\zihao{5},
}

% ---- 页面尺寸 ----
\usepackage[
  a4paper,
  top=19mm, bottom=25mm,
  left=15mm, right=15mm,
  columnsep=5mm
]{geometry}

% ---- 字体 ----
\usepackage{fontspec}
\setmainfont{Times New Roman}
\setsansfont{Arial}
\setmonofont{Courier New}

% ---- 行距 ----
\usepackage{setspace}
\setstretch{1.1}

% ---- 图表 ----
\usepackage{graphicx}
\usepackage{booktabs}
\usepackage{multirow}
\usepackage{float}
\usepackage{subfigure}

% ---- 数学 ----
\usepackage{amsmath, amssymb, amsthm}
\newtheorem{theorem}{定理}
\newtheorem{lemma}{引理}
\newtheorem{definition}{定义}

% ---- 算法 ----
\usepackage[ruled, linesnumbered]{algorithm2e}
\SetAlFnt{\small}

% ---- 超链接 ----
\usepackage[
  colorlinks=true,
  linkcolor=black,
  citecolor=blue,
  urlcolor=blue
]{hyperref}

% ---- 参考文献 ----
\usepackage[numbers, sort&compress]{natbib}
\bibliographystyle{unsrtnat}

% ---- 页眉页脚：仅保留页脚页码 ----
\usepackage{fancyhdr}
\pagestyle{fancy}
\fancyhf{}
\fancyfoot[C]{\small\thepage}
\renewcommand{\headrulewidth}{0pt}
\renewcommand{\footrulewidth}{0pt}

% ---- 标题信息（按需修改）----
\title{
  \vspace{-1em}
  {\heiti\zihao{2} 论文标题}\\[4pt]
  {\zihao{-4} Paper Title in English}
}
\author{
  \zihao{-4}
  作者一$^{1}$\quad 作者二$^{2}$\\[2pt]
  \zihao{5}
  $^{1}$单位一，城市 邮编\\
  $^{2}$单位二，城市 邮编\\[2pt]
  \texttt{\zihao{5} email@example.com}
}
\date{}

% ============================================================
\begin{document}

\maketitle
\thispagestyle{fancy}

% ============================================================
% 本节介绍模版整体结构
% ============================================================
\section{模版说明}

本模版适用于双栏中文学术论文写作，风格参考 CCF-A 会议与 IEEE 期刊排版规范。
文档类为 \texttt{article}，使用 \texttt{ctex} 宏包提供中文支持，
编译方式为 XeLaTeX。

% ============================================================
\section{文档结构}

论文正文结构不固定，按需使用以下命令划分层次：

\begin{itemize}
  \item \verb|\section{标题}| — 一级节，黑体小四加粗
  \item \verb|\subsection{标题}| — 二级节，黑体五号加粗
  \item \verb|\subsubsection{标题}| — 三级节，楷书五号
\end{itemize}

% ============================================================
\section{字体与字号}

正文使用 10pt Times New Roman（英文）与宋体（中文）。
用 \verb|\zihao{n}| 手动切换字号，常用对照如下：

\begin{itemize}
  \item \verb|\zihao{2}| — 二号，标题
  \item \verb|\zihao{-4}| — 小四，作者行
  \item \verb|\zihao{5}| — 五号，图表说明、脚注
\end{itemize}

用 \verb|\heiti|、\verb|\songti|、\verb|\kaishu| 切换黑体、宋体、楷书。
英文加粗用 \verb|\textbf{}|，斜体用 \verb|\textit{}|。

% ============================================================
\section{数学公式}

行内公式用 \verb|$...$|，例如损失函数 $\mathcal{L} = \frac{1}{n}\sum_i \ell_i$。

独立编号公式用 \texttt{equation} 环境：
\begin{equation}
  \mathcal{L} = \frac{1}{n}\sum_{i=1}^{n} \ell\!\left(f(x_i),\, y_i\right)
  \label{eq:example}
\end{equation}

用 \verb|\eqref{eq:example}| 引用公式，得到式~\eqref{eq:example}。
多行对齐公式使用 \texttt{align} 环境。

% ============================================================
\section{图与表}

\subsection{插图}

用 \texttt{figure} 环境插图，位置参数推荐 \verb|[t]|（当前栏顶部）。
跨双栏大图用 \verb|figure*| 环境。

\begin{figure}[t]
  \centering
  % 取消注释并替换路径：
  % \includegraphics[width=\linewidth]{figs/example.pdf}
  \fbox{\rule{0pt}{2.5cm}\rule{\linewidth}{0pt}}
  \caption{图题写在图下方，五号字。}
  \label{fig:example}
\end{figure}

用 \verb|\ref{fig:example}| 引用，得到图~\ref{fig:example}。

\subsection{表格}

表题写在表格上方。推荐使用 \texttt{booktabs} 的三线表风格。

\begin{table}[t]
  \centering
  \caption{表题写在表格上方}
  \label{tab:example}
  \zihao{5}
  \begin{tabular}{lcc}
    \toprule
    方法 & 指标A & 指标B \\
    \midrule
    基线方法 & 72.3 & 68.1 \\
    \textbf{本文方法} & \textbf{78.9} & \textbf{75.6} \\
    \bottomrule
  \end{tabular}
\end{table}

% ============================================================
\section{算法}

用 \texttt{algorithm2e} 环境描述算法，位置参数用 \verb|[t]|。

\begin{algorithm}[t]
\caption{算法示例}
\label{alg:example}
\KwIn{输入数据 $\mathcal{D}$，学习率 $\eta$}
\KwOut{模型参数 $\theta$}
初始化 $\theta$\;
\For{每个训练轮次}{
  采样小批量 $\mathcal{B} \subset \mathcal{D}$\;
  $\theta \leftarrow \theta - \eta \nabla_\theta \mathcal{L}(\mathcal{B};\theta)$\;
}
\Return $\theta$\;
\end{algorithm}

% ============================================================
\section{参考文献}

参考文献统一管理在 \texttt{refs.bib} 文件中，正文用
\verb|\cite{key}| 引用，编译后自动生成编号。

示例引用：深度学习综述见~\cite{lecun2015deep}，
Transformer 架构见~\cite{vaswani2017attention}，
ResNet 见~\cite{he2016deep}。

引用风格为按出现顺序排列的数字格式（\texttt{unsrtnat}），
若需要作者-年份格式，将 \verb|\bibliographystyle| 改为
\texttt{plainnat} 并去掉 \texttt{natbib} 的 \texttt{numbers} 选项。

% ---- 参考文献 ----
\bibliography{refs}

\end{document}
